% Drop-in result and claim revisions based on the two reviewed full studies.
% This fragment does not replace edits made in the other manuscript chat.
% Requires the manuscript's \paragraphB command; all prose paragraphs are unwrapped.
% Full raw replay remains pending. Do not claim that it has passed until its report is returned.

\paragraphB{Timing interference under equal competing service.} We compare release policies using identical competing programs and completed service. Across 16,384 fresh paired trials per workload and arm, learned timing increases teleportation measurement errors by 2.02 percentage points relative to random timing, with a simultaneous 95\% interval of [1.29,2.75] points, and by 2.66 points relative to even spacing, with interval [1.93,3.39]. For the plus-state workload, the increase relative to random timing is 0.52 points with interval [-0.17,1.22], whereas the increase relative to even spacing is 1.56 points with interval [0.89,2.22]. The estimates are conditional on the frozen ensemble of previously calibrated policies and the declared public timing window.

\paragraphB{What calibration adds.} The fixed-phase comparison limits the advantage attributable to calibration. For the plus-state workload, learned error minus the largest fixed-grid error has a simultaneous envelope of [-1.47,-0.04] percentage points. For teleportation, the envelope is [-0.79,0.33] points. The fixed-grid maximum is an oracle diagnostic selected from the evaluated phases; its phase is not assumed to be available to an attacker beforehand. These findings establish an advantage over random or evenly spaced interference in the stated comparisons, without establishing that calibration outperforms the strongest fixed phase.

\paragraphB{Application-level effects.} The application comparison uses 8,192 independent paired streams per condition, each containing 12 offered quantum requests. In the short-memory configuration with a 10\,$\mu$s host instruction cost and unoptimized programs, selected interference increases mean infidelity over random interference by 0.421 percentage points. Its simultaneous 95\% interval, [-0.0033,0.846] points, includes zero. At the source-default host cost and memory parameters, the point estimate is negligible. The tested delayed-preparation and rotation-merging transformation yields identical observed infidelity across the active interference arms. This boundary concerns the evaluated rotation application; it does not establish protection for arbitrary interactive protocols.

\paragraphB{Quality and useful service.} Under selected interference in the short-memory, 10\,$\mu$s unoptimized configuration, numerical earliest-deadline-first scheduling improves correct-and-timely completion over first-come-first-served scheduling by 4.18 percentage points, with a simultaneous interval of [3.54,4.81]. The live-state priority policy performs worse than numerical earliest-deadline-first scheduling by 1.92 points, with interval [1.37,2.47]. Critical-section exclusion reduces mean infidelity relative to first-come-first-served scheduling by 9.29 points, with interval [8.87,9.71], but also reduces correct-and-timely completion by 13.45 points, with interval [12.39,14.51]. Its observed late-result rate is 25.92\%, compared with 1.79\% under first-come-first-served scheduling and zero under numerical earliest-deadline-first scheduling. These results show why quality must be assessed together with service outcomes.

\paragraphB{Causal controls and sensitivity.} Separate diagnostic studies comprise 512 independent causal-control streams and 256 sensitivity streams. For unoptimized first-come-first-served scheduling, the selected-minus-random mean-infidelity estimates are zero with ideal operations or gate noise alone, +0.391 percentage points with memory noise alone, and +0.359 points with combined noise. This pattern is consistent with a contribution from memory exposure. However, every interval in the predeclared families of 140 causal and 450 sensitivity contrasts includes zero. The optimized infidelity contrasts have zero point estimates throughout both studies. We retain these findings as diagnostic evidence and do not interpret them as precise confirmation of robustness or population-wide absence of an effect.

\paragraphB{Statistical unit and scope.} Timing inference uses a fresh paired trial; application inference uses an entire independent stream. We fix the sample sizes and comparison families before evaluation, and report two-sided empirical Bernstein intervals with a union bound over 36 timing comparisons or six primary application comparisons. For paired differences $X_i\in[a,b]$, the radius is $\sqrt{2S_n^2\log(4K/\alpha)/n}+7(b-a)\log(4K/\alpha)/(3(n-1))$, where $S_n^2$ is the sample variance, $K$ is the declared family size, and $\alpha=0.05$. Primary differences use $[a,b]=[-1,1]$; diagnostic differences of differences use $[-2,2]$. We count all offered requests and preserve wrong-and-timely, late, unfinished, and wrong-and-late outcomes. Source-default parameters and deliberately changed sensitivity parameters are reported separately. Exact software versions and source fingerprints belong in the replication artifact; all essential experimental assumptions and analysis choices remain in the main paper.

% Suggested contribution wording. Use these to replace stronger unsupported claims.
% 1. We quantify when the release timing of equal-service classical work changes quantum measurement outcomes on a shared simulated network node.
% 2. We characterize how the observed effect depends on workload, timing information, memory assumptions, and applicable program transformations.
% 3. We measure the tradeoff between quantum quality and useful service across established scheduling policies and critical-section exclusion, including cases where the proposed priority policy underperforms numerical EDF.
% Claims to remove: calibration universally beats fixed timing; the live-state priority policy beats EDF; the attack generally survives optimization; complete upstream compiler/TUF methods were reproduced; deployed permissions or device timings were validated.
